\documentclass[12pt]{article}
\usepackage[utf8]{inputenc}
\usepackage[T1]{fontenc}
\usepackage{hyperref}
\title{Ararandew\'ara Language Resources for K-12 Curriculum Development}
\author{Compiled by Claude Sonnet 5.5 (claude-sonnet-5-5) for Andr\'as Kornai}
\date{}
\begin{document}
\maketitle

\section{Introduction}
Ararandew\'ara (ISO 639-3 xaj) is an extinct Tupian language, attested in 1913--14 (Lang 1914) and formerly spoken at the headwaters of the Moju River, Par\'a, Brazil. A source on Tupian languages describes it as closely related to, but possibly distinct from, Amanay\'e. It should not be confused with Karo, the language of the Arara people of Rond\^onia. The sources found are very thin.

\section{Primary Educational Resources}
None found.

\section{Linguistic and Reference Materials}
\begin{itemize}
\item Ararandew\'ara language article: \url{https://en.wikipedia.org/wiki/Ararandew%C3%A1ra_language}
\item Tupi language overview: \url{https://en.wikipedia.org/wiki/Tupi_language}
\item Museu Goeldi repository record, returned by the search; its relevance was not verified: \url{https://repositorio.museu-goeldi.br/handle/mgoeldi/1219}
\end{itemize}

\section{Religious and Community Resources}
None found.

\section{Digital and Audio Resources}
None found.

\section{Curriculum Development Context}
There is almost no linguistic documentation, and I found no descendant community or revitalization programme. A K-12 programme could at most present the historical record and place the language within the Tupian and Amanay\'e context. The search surfaced no usable vocabulary or grammar.
\end{document}